\chapter{现代大模型研究方向简介}\label{现代大模型的研究方向简介}

前十章沿着训练与使用一个学习助手的过程，介绍了数据、模型、后训练、推理、检索、内部分析与多模态建模。本章把视野扩展到整个大模型研究领域，按研究问题列出主要方向。每个方向简要交代问题从何而来、为什么值得研究，以及已经形成的方法与代表成果，供读者选择下一阶段的阅读入口。

这些方向经常相互连接。例如，长任务智能体同时涉及推理、工具使用、记忆、强化学习与系统效率；多模态助手同时需要感知、跨模态对齐与交互评估。读者可以先选择一个感兴趣的问题，再沿着它依赖的数据、算法和系统环节继续深入。

\section{数据与预训练 能力从哪里获得}

\begin{itemize}
    \item \textbf{数据质量与数据配比。}训练语料覆盖网页、书籍、代码、数学与多语言文本，不同来源的价值与重复程度差异很大。研究目标是在给定训练预算下选择更有效的内容。去重、质量过滤、领域混合与数据消融已成为预训练的重要组成，FineWeb、Dolma 等开放语料项目提供了可研究的数据处理流程。
    \item \textbf{合成数据与数据自举。}高质量示范和困难任务标签获取成本较高，研究者让模型生成指令、解题过程和交互轨迹，再通过规则、执行结果或模型评分筛选。Self-Instruct、拒绝采样和可验证任务生成形成了常见路线，代码执行与数学验证使部分合成数据能够自动检查。
    \item \textbf{规模规律与计算最优训练。}参数数量、数据量和训练步数共同消耗资源，研究者希望预测如何分配预算。Scaling laws 与 Chinchilla 等工作建立了损失、模型规模和数据规模之间的经验关系，为训练规模选择与小规模预测实验提供了工具。
    \item \textbf{分词、语言覆盖与低资源语言。}同一句话在不同分词器中可能占用不同数量的 token，影响训练效率与使用成本。字节级表示、多语言词表、跨语言迁移和语言均衡采样，推动模型覆盖更多语言、方言与混合语言输入。
    \item \textbf{训练目标与课程学习。}下一 token 预测、去噪重建和多 token 预测提供不同的学习信号；样本顺序与难度也影响优化过程。研究者通过混合目标、难度调度和阶段式数据配比，探索更高的数据利用率与能力形成效率。
\end{itemize}

\section{模型架构 怎样组织参数与记忆}

\begin{itemize}
    \item \textbf{稠密模型与混合专家。}扩大模型容量会增加计算开销。混合专家（Mixture of Experts，MoE）通过路由器为每个 token 选择部分专家，使总参数容量与单次激活计算量分开增长。Switch Transformer、Mixtral 和 DeepSeek 系列推动了稀疏专家的规模化应用，路由均衡、专家分工与通信效率是主要研究问题。
    \item \textbf{注意力与历史压缩。}长序列带来大量历史读取和缓存占用。分组查询注意力、多头潜在注意力、滑动窗口与稀疏注意力分别从共享、压缩和访问范围降低成本。相关成果已经进入主流模型架构，研究继续围绕检索精度、记忆容量与硬件效率展开。
    \item \textbf{线性注意力与状态空间模型。}持续生成与资源受限部署需要控制历史状态的增长。线性注意力和状态空间模型把历史信息写入可递推更新的状态，Mamba、DeltaNet 与 Kimi Linear 等工作形成了代表路线。混合架构进一步组合状态记忆与逐位置注意力。
    \item \textbf{位置表示与长上下文。}模型需要区分顺序、距离以及远处证据。旋转位置嵌入及其扩展、长序列继续训练和分层记忆，使模型能够处理更长文档。新的研究重点包括跨段整合、远距离精确检索和长度增长后的稳定表现。
    \item \textbf{替代生成范式。}逐 token 自回归生成形成串行依赖，扩散语言模型与掩码迭代生成则研究同时修改多个位置。离散扩散与块级生成已经用于文本建模实验，研究目标包括并行解码、可控修改与生成质量之间的协调。
\end{itemize}

\section{训练系统 怎样把算法变成可运行的大规模计算}

\begin{itemize}
    \item \textbf{分布式并行与通信。}单张加速卡难以容纳大型模型及其优化器状态。数据并行、张量并行、流水线并行、专家并行和序列并行分别切分不同计算对象；ZeRO、FSDP 与 Megatron 系列形成了代表性系统。研究关注显存分配、通信重叠、负载均衡和设备故障恢复。
    \item \textbf{低精度训练与优化稳定性。}降低数值精度可以减少存储和带宽需求，同时改变误差积累方式。混合精度、损失缩放和 FP8 训练已形成实践路线，优化器、梯度异常检测与训练恢复机制共同影响长时间训练的稳定性。
    \item \textbf{算子、编译器与硬件协同。}同一算法的运行速度受访存、算子融合与设备布局影响。FlashAttention 展示了通过组织计算与数据访问提高效率的路径，自动内核生成、图编译与硬件感知调度继续降低模型训练成本。
\end{itemize}

\section{后训练与对齐 怎样把基础能力变成有用行为}

\begin{itemize}
    \item \textbf{指令微调与行为数据。}预训练模型需要学习任务要求、对话格式和输出结构。监督微调、指令混合与多轮示范已经成为助手训练的基础环节，数据覆盖、难度配比与能力保持决定后训练的收益。
    \item \textbf{偏好学习与奖励建模。}答案质量同时涉及正确性、清晰度、风格和用户需求。RLHF、DPO、IPO 等方法把比较标签转化为策略更新，研究重点包括偏好噪声、奖励泛化、多目标冲突和不同用户的偏好差异。
    \item \textbf{可验证奖励与推理强化学习。}数学答案、程序测试和工具任务提供了自动反馈。GRPO、PPO 及相关变体推动模型通过尝试学习，高熵位置选择、题目难度调度、奖励设计与探索维持成为活跃问题。
    \item \textbf{过程监督与信用分配。}最终答案反馈需要传到中间决策。过程奖励、步骤验证、优势估计与轨迹比较尝试识别哪些步骤推动成功，相关研究已经形成步骤标注数据和过程验证器，并继续探索自动化监督。
    \item \textbf{蒸馏与能力迁移。}强模型的计算成本推动知识向小模型迁移。回答蒸馏、长推理蒸馏和同策略蒸馏分别利用示范序列或学生轨迹上的教师反馈，形成了面向低成本推理与特定任务部署的训练路线。
\end{itemize}

\section{推理与搜索 怎样分配一次求解的计算}

\begin{itemize}
    \item \textbf{推理时扩展。}困难问题通常需要更多尝试与检查。思维链、自一致性、多候选择优和自适应推理预算，把固定模型的计算分配到不同解题路径。研究关注任务难度、采样数量、路径长度与成功率的关系。
    \item \textbf{搜索与验证器。}多步任务包含分支选择，验证器可以评估中间状态或最终候选。树搜索、束搜索与验证器引导生成，将语言模型连接到显式搜索过程；数学检查器和程序测试提供了可执行的候选评估。
    \item \textbf{潜在空间推理与自适应计算。}把全部中间计算写成文字会产生较长序列，研究者探索在隐藏状态中迭代、复用层或动态分配深度。已有工作给出了连续表示推理与循环计算的实验方案，目标是提高单位计算的求解能力。
    \item \textbf{神经符号结合与形式化证明。}语言模型擅长提出候选步骤，符号系统擅长精确执行与规则检查。程序辅助推理、定理证明器协作与形式化数学，将候选生成和机器验证连接起来，形成可搜索、可检查的求解流程。
\end{itemize}

\section{工具与智能体 怎样在环境中完成任务}

\begin{itemize}
    \item \textbf{工具选择与结构化调用。}外部信息和精确执行需要搜索、数据库与代码环境。ReAct、Toolformer 和函数调用训练建立了从生成调用到读取结果的基本流程，研究重点转向参数可靠性、多工具组合和执行错误恢复。
    \item \textbf{长任务规划与状态管理。}软件修改、资料研究和业务流程需要跨越许多步骤。任务分解、计划更新、检查点与结构化状态让系统持续跟踪进度，长期交互数据与任务级评测推动模型从单轮回答走向完整任务执行。
    \item \textbf{浏览器、桌面与软件工程智能体。}许多应用通过图形界面、终端或代码仓库提供操作入口。界面定位、程序修改和执行反馈构成训练环境，WebArena、OSWorld 与 SWE-bench 等基准分别提供网页、桌面和代码任务的评估入口。
    \item \textbf{多智能体协作。}复杂任务可以按检索、求解、执行和审查等角色分工。研究探索通信协议、角色配置、共享记忆和协作奖励，已有系统实现了多角色任务流水线。通信开销、重复工作和错误传播是优化协作结构的核心问题。
    \item \textbf{环境学习与自我改进。}智能体可以从工具反馈和任务记录中积累经验。在线强化学习、自动课程、失败回放和经验检索，形成持续改进策略与执行流程的路线，研究关注反馈质量与经验在新任务中的复用。
\end{itemize}

\section{检索、记忆与持续学习 怎样获取和更新知识}

\begin{itemize}
    \item \textbf{检索增强与证据组织。}企业资料和新事实需要随时更新。混合检索、重排、查询改写、GraphRAG 与多轮检索已经形成较完整的应用方法，研究进一步关注跨文档证据组合、时效冲突与引用质量。
    \item \textbf{长期记忆与个性化。}跨会话助手需要保存用户偏好、历史决策和任务经验。外部记忆库、事件摘要与记忆检索构成基本方案，研究问题包括什么时候写入、如何更新、哪些信息应遗忘，以及如何保持用户控制。
    \item \textbf{持续学习与领域适配。}新知识、新领域和新任务不断出现，重复全量训练成本较高。领域继续预训练、参数高效微调与经验回放提供了增量适配方式，知识获取、遗忘控制和任务迁移构成主要评估目标。
    \item \textbf{模型编辑与机器遗忘。}特定事实修正和数据删除提出了更局部的更新需求。ROME、MEMIT 等工作探索参数编辑，机器遗忘研究移除指定数据影响的训练方法。定位更新范围、评估泛化以及测量相邻能力变化，是这些方向的重要问题。
\end{itemize}

\section{多模态与具身智能 从语言走向感知和行动}

\begin{itemize}
    \item \textbf{图像理解与文档理解。}图片、图表和页面布局包含文字之外的信息。视觉编码器、跨模态连接模块与视觉语言指令训练，推动了 LLaVA、Qwen-VL 等模型路线，支持视觉问答、图表解释和文档阅读。
    \item \textbf{视频理解与时序推理。}动作先后、事件持续时间和跨帧关系需要时间建模。帧选择、视频压缩与长序列融合形成了主要方法，相关模型已支持视频问答、片段定位与事件摘要，持续视频输入推动流式记忆研究。
    \item \textbf{语音与实时交互。}自然对话同时包含语义、语气、节奏和轮次切换。语音 token、端到端语音建模与流式生成推动口语助手发展，研究关注低延迟、打断处理、情绪表达和多说话人场景。
    \item \textbf{统一理解与生成。}用户既希望理解图片和声音，也希望生成或编辑它们。自回归与扩散结合、统一多模态表示和联合训练，使理解与生成进入共享模型；主要问题是不同模态目标的协同和生成过程的精确控制。
    \item \textbf{世界模型与空间智能。}规划需要预估行动怎样改变环境。视频预测、三维表示与交互动态建模尝试学习环境状态转移，生成式模拟器提供了可观察的研究成果，物理一致性、长时预测和行动条件控制是重要课题。
    \item \textbf{机器人与视觉语言动作模型。}机器人需要把指令与视觉观察转成连续动作。视觉语言动作（Vision-Language-Action，VLA）模型、机器人轨迹预训练与仿真迁移形成代表路线，RT-2、OpenVLA 等工作连接了语言知识与机器人控制。
\end{itemize}

\section{部署与效率 怎样降低一次服务的成本}

\begin{itemize}
    \item \textbf{量化、剪枝与小模型。}终端设备与高并发服务要求较小的存储和计算量。GPTQ、AWQ、结构化剪枝和蒸馏提供模型压缩方法，研究重点是不同层的精度分配、任务质量保持与实际硬件加速。
    \item \textbf{推测解码与并行生成。}自回归解码中，主模型逐步调用形成延迟。草稿模型先提出多个 token，主模型批量验证；推测解码与自推测方法已经进入推理引擎，研究继续优化草稿成本、接受率与批处理协同。
    \item \textbf{缓存、批处理与调度。}在线请求的长度、到达时间和工具等待各不相同。分页缓存、前缀复用、连续批处理和预填充与解码分离，提升了显存利用与服务吞吐。研究同时关注单用户延迟、总处理量和请求公平性。
    \item \textbf{端侧、云端与能源效率。}本地运行提供低延迟与数据就地处理，云端提供更大的计算资源。模型路由、大小模型协同、缓存下沉与能耗测量，帮助系统按任务难度和设备能力分配计算。
\end{itemize}

\section{评估、可靠性与安全 怎样衡量模型的实际表现}

\begin{itemize}
    \item \textbf{评测设计与数据污染。}公开题目可能进入训练集，单一分数也难覆盖复杂行为。动态题集、时间切分、私有测试和污染检测推动评估改进，任务级基准进一步测量完整流程的成功率与成本。
    \item \textbf{自动评审与人类评估。}开放式回答难以通过字符串匹配评分。模型评审、成对比较与人工标注扩大了评测范围，研究关注评审器的位置偏好、风格偏好、一致性和与真实用户判断的关系。
    \item \textbf{事实性、不确定性与校准。}模型需要区分已有证据、推断结果与信息缺口。检索核验、置信度估计、选择性回答和校准方法提供了改进路线，目标是让回答强度与实际正确概率相匹配。
    \item \textbf{对抗鲁棒性与提示注入。}复杂输入、外部文档与工具返回会影响模型决策。对抗训练、输入来源标记、权限隔离和调用验证形成系统研究方向，评测覆盖越狱、跨文档注入与多轮攻击等情形。
    \item \textbf{隐私、数据归属与内容溯源。}训练和部署涉及私人数据、版权材料与生成内容传播。成员推断评估、隐私保护训练、机器遗忘、水印与来源记录，分别研究数据暴露、删除和生成内容追踪问题。
    \item \textbf{公平性、价值偏好与社会影响。}不同群体、语言和使用情境对回答有不同要求。分群评估、多元偏好建模与参与式数据收集，帮助研究者识别表现差异，分析模型部署对劳动、教育和信息传播的影响。
\end{itemize}

\section{机制、理论与科学应用 怎样理解能力并拓展用途}

\begin{itemize}
    \item \textbf{机制可解释性与表示结构。}模型内部如何编码概念、路由信息和形成答案，是理解行为的重要问题。探针、激活修补、稀疏自编码器与回路追踪已经提供可操作的实验工具，研究继续向自动化分析与任务级机制验证发展。
    \item \textbf{上下文学习、泛化与记忆。}模型为什么能从少量示例学习、何时依赖记忆、何时能组合已有知识，涉及学习机制本身。受控合成任务、归纳头分析与训练动力学实验，为区分数据记忆、算法学习和分布迁移提供了研究手段。
    \item \textbf{训练动力学与能力形成。}能力在训练的哪些阶段出现，数据顺序、优化器与模型结构怎样影响学习，是连接理论和实践的问题。检查点分析、损失预测与能力消融已经形成研究方法，有助于解释训练预算和行为变化之间的关系。
    \item \textbf{代码、数学与科学发现。}程序执行、形式化证明和科学模拟器提供清晰反馈，使模型能够参与更复杂的搜索。代码助手、自动证明与实验方案生成形成了应用成果，研究关注从文本建议走向可执行、可验证的研究产物。
    \item \textbf{专业领域与人机协作。}教育、医疗、法律和工程具有各自的知识结构、工作流和质量要求。领域适配、证据检索与专业工具集成推动专用助手发展，用户研究和流程级评估帮助判断模型怎样改善真实工作。
\end{itemize}

\section{怎样选择自己的深入方向}

选择方向时，可以先确定自己最想解释或改善的现象。对数学解题感兴趣，可以从奖励设计、推理搜索和高熵位置训练切入；对系统性能感兴趣，可以选择缓存、量化或解码；对知识问答感兴趣，可以研究检索、证据组织和更新；对模型内部计算感兴趣，可以沿第九章的探针与干预实验继续深入。

接着确定可获得的资源。已有模型和少量数据适合开展评测、检索、提示、工具调用与小规模内部分析；具备单卡训练条件时，可以研究参数高效微调、小模型蒸馏和受控训练实验；分布式资源支持更大规模的预训练、强化学习和架构比较。清晰的问题、可复现的基线与能够解释差异的实验，是不同资源条件下共同的研究起点。

阅读可以按``综述建立地图、代表论文理解方法、实现代码核对细节、后续工作追踪问题''的顺序推进。为一个方向整理背景、主要假设、数据来源、评价指标和代表成果，再选一个变量做对照实验，读者就能从了解方向逐步进入研究实践。

\textbf{进一步阅读入口}
\begin{enumerate}
    \item Zhao 等，\href{https://arxiv.org/abs/2303.18223}{A Survey of Large Language Models}。覆盖预训练、适配、能力与评估，适合建立总体研究地图。
    \item DeepSeek-AI，\href{https://arxiv.org/abs/2412.19437}{DeepSeek-V3 Technical Report}。结合模型架构与训练系统，阅读大规模模型怎样组织完整技术方案。
    \item \href{https://arxiv.org/abs/2503.21460}{Large Language Model Agent: A Survey on Methodology, Applications and Challenges}。沿智能体结构、工具交互和评估继续阅读。
    \item Yin 等，\href{https://arxiv.org/abs/2306.13549}{A Survey on Multimodal Large Language Models}。了解多模态模型的连接方式、训练数据与任务。
    \item 前文第五至九章的相关资料分别提供强化学习与偏好优化、知识蒸馏、推理效率、RAG 与可解释性的专题入口，可按本章的问题分类交叉阅读。
\end{enumerate}
